\documentclass[10pt]{article}

\usepackage[utf8]{inputenc}

\usepackage[T1]{fontenc}

\usepackage{amsmath}

\usepackage{amsfonts}

\usepackage{amssymb}

\usepackage[version=4]{mhchem}

\usepackage{stmaryrd}

\usepackage{CJKutf8}



\begin{document}

С помощью Ц. ф. выражается общее решение линейного уравнения



\[

z^{2} W^{\prime \prime}+(1-2 \alpha) z W^{\prime}+\left[\beta^{2} \gamma^{2} z^{2 \gamma}-\left(p^{2} \gamma^{2}-\alpha^{2}\right)\right] W=0

\]



при произвольных $\alpha, \beta, \gamma$. Именно:



\[

W=c_{1} z^{\alpha} Z_{p}^{(1)}\left(\beta z^{\gamma}\right)+c_{2} z^{\alpha} Z_{p}^{(2)}\left(\beta z^{\gamma}\right) .

\]



Среди уравнений (2) содержится уравнение вида



\[

W^{\prime}-z W=0

\]



к-рое порождает так наз. фу н кци и Э й р и.



В приложениях часто встречаются м о д и ф и ц и р о в а н ны е Ц. ф. (Ц. ф. мнимого аргумента), к-рые определяются соотношением:



\[

I_{p}(x)=e^{-\frac{p \pi i}{2}} J_{p}(i x)=\sum_{k=0}^{\infty} \frac{1}{k ! \Gamma(p+k+1)}\left(\frac{x}{2}\right)^{p+2 k}

\]



и удовлетворяют уравнению вида (2) при $\alpha=0, \beta=i, \gamma=1$ :



\[

U^{\prime \prime}+\frac{1}{x} U^{\prime}-\left(1+\frac{p^{2}}{x^{2}}\right) U=0

\]



в качестве второго решения берут здесь или $I_{-p}(x)$ (при нецелом $p$ ), или функцию Макдональда



\[

K_{p}(x)=\frac{\pi i}{2} e^{\frac{p \pi i}{2}} H_{p}^{(1)}(i x) \text {. }

\]



Важность этих функций обусловлена тем, что они представляют собой положительные решения уравнения (3) при $x>0$ с экспоненциальным поведением на бесконечности, асимптотич. представления имеют для них вид:



\[

I_{p}(x) \sim \frac{1}{\sqrt{2 \pi x}} e^{x} ; \quad K_{p}(x) \sim \sqrt{\frac{\pi}{2 x}} e^{-x} .

\]



В нек-рых задачах встречаются еще Ц. Ф. аргумента



\[

z=x \sqrt{-i}=e^{3 i \pi / 4} x

\]



Для действительных и мнимых частей этих функций введены специальные обозначения ( $p$ - действительно):



\[

\begin{gathered}

I_{p}(x \sqrt{-i})=\operatorname{ber}_{p} x+i \operatorname{bei}_{p} x, \\

K_{p}(x \sqrt{-i})=\operatorname{ker}_{p} x+i \operatorname{kei}_{p} x, \\

H_{p}(x \sqrt{-i})=\operatorname{her}_{p} x+i \operatorname{hei}_{p} x .

\end{gathered}

\]



Лит.: [1] Лаврентьев М. А., Шабат Б. В., Методы теории функций комплексного переменного, 3 изд., М., 1965; [2] Л е бе д е в Н. Н., Специальные функции и их приложения, 2 изд., М.- Л., 1963; [3] В а тсо н Г. Н., Теория бесселевых функций, пер. с англ., ч. 1-2, М., 1949; [4] Я нке Е., Эм де Ф., Л е ш Ф., Специальные функции. Формулы, графики, таблицы, 2 изд., пер. с нем., М., 1968.



П. И. Лизоркин.



ЧИЛИНДРИЧЕСКОЕ МНОХЕСТВО - см. Случайный процесс, Случайное поле.



ЦИРКУЛЯРНЫЙ ИНТЕТРАЛ - см. Эллиптический интегpar.



ЦИРКУЛЯ\begin{CJK}{UTF8}{mj}山\end{CJK}ИЯ векторного поля $\boldsymbol{a}(\boldsymbol{r})$ вдоль замкнутой кривой $L$ - интеграл вида:



\[

\oint_{L} a d r

\]



в координатной форме Ц. равна



\[

\int_{L}\left(a_{x} d x+a_{y} d y+a_{z} d z\right)

\]



Работа, совершаемая силами силового поля $\boldsymbol{a}(\boldsymbol{r})$ при перемешении пробного тела (единичной массы, заряда и т. д.) вдоль кривой $L$, равна Ц. поля вдоль $L$.



См. также Векторный анализ.



В. И. Битюцков. ЧАНДРАСЕКАРА - МЮНША УРАВНЕНИЕ - уравнение для плотности распределения вероятностей $f(u ; x)$ наблюдаемой яркости $\xi$ звездного скопления, находящегося на расстоянии $x$ от Земли, $\int_{0}^{\infty} f(u ; x) d u=1$. В теории Чандрасекара - Мюнша яркость $\xi$, наблюдаемая с Земли в заданном направлении, вследствие случайного распределения туманностей во Вселенной со случайным коэффициентом прозрачности $q<1$ рассматривается как случайный разрывный марковский процесс $\{\xi(x), x \in(0, \infty)\}$. Тогда



\[

f(u ; x) d u=\mathrm{P}\{\xi(x) \in[u, u+d u]\} .

\]



При известной плотности $f(u ; x-0)$ в точках $x-h, h \rightarrow+0$, условная вероятность события $\{\xi(x+0) \in[u, u+d u]\}$ при условии, что в точке $x$ свет проходит через туманность с коэффициентом прозрачности $q$, равна $f(u / q ; x) d u / q$. Полагая, что моменты встречи луча света с туманностью составляют пуассоновский поток событий с интенсивностью $\alpha$, статистически независимый от $q$ и $\psi(q)$ - плотность распределения $q, \int_{0}^{1} \psi(q) d q=1$, получим, что процесс $\{\xi(x)\}$ - марковский. Из уравнения Колмогорова - Чепме-



\section{2 ЦИЛИНДРИЧЕСКИЕ}

на вытекает:



\[

\begin{gathered}

f(u ; x+h)=f(u-\beta h ; x)(1-\alpha h)+ \\

+\alpha h \int_{0}^{1} d q \psi(q) f(u / q ; x) / q+o(h), \quad h \rightarrow+0,

\end{gathered}

\]



где $\beta$ - коэффициент детерминированного угасания яркости на единице длины. Из (1) следует Ч.- М. у.



\[

\begin{gathered}

\frac{\partial}{\partial x} f(u ; x)+\beta \frac{\partial}{\partial u} f(u ; x)= \\

=\alpha\left[\int_{0}^{1} f(u / q ; x) \psi(q) d q / q-f(u ; x)\right] .

\end{gathered}

\]



Так как $\xi(0)=0$, то в качестве начального условия для этого уравнения нужно взять $f(u ; 0)=\delta(u)$ и нижний предел в интеграле в правой части (2) может быть заменен на $u / \beta x$.



Лит.: [1] Б а р у ч а-Р и д А. Т., Элементы теории марковских процессов и их приложения, пер. с англ., М., 1969; [2] C ha ndrase kar S., Munch G., "Astrophys. J.», 1950, v. 112, № 3, p. $380-92$, 393-98; 1951, v. 114, № 1, p. 110-22; 1952, v. 115, № 1, p. $94-102,103-23$



Ю. П. Вирченко.



ЧАПЛЬIТИА ЗАДАЧА о д в и Ж ни и т ве рдого те л а в жидкости - задача, решение которой найдено С. А. Чаплыгиным в 1903 (см. [1]). Ч. з. отвечает случаю, когда живая сила имеет вид:



\[

T=\frac{a}{2}\left(P_{1}^{2}+P_{2}^{2}+2 P_{3}^{2}+(b+c) R_{1}^{2}+(b-c) R_{2}^{2}+b R_{3}^{2}\right),

\]





\end{document}